%=====================================================================
% PART II OPENER
%=====================================================================
\thispagestyle{plain}
\structuralfigures{P2.}

\begin{center}
\begin{tcolorbox}[enhanced, width=0.92\textwidth, colback=secondary!5,
  colframe=secondary, boxrule=0pt, leftrule=5pt, arc=3pt, top=8pt, bottom=8pt]
\large\itshape\color{secondary!80!black}
An algorithm that calls itself is easy to write and hard to cost. Part~II is the machinery for costing one --- three ways to solve a recurrence, a theorem that solves most of them at sight, and the two sorting algorithms that every later chapter uses as a reference point.
\end{tcolorbox}
\end{center}

\vspace{6pt}

\dsfigh{%
\begin{tikzpicture}[
  cbox/.style={rounded corners=3pt, draw=secondary, line width=1pt, fill=secondary!7,
               text=secondary!80!black, font=\small\bfseries, align=center,
               minimum width=3.0cm, minimum height=1.0cm},
  hbox/.style={rounded corners=3pt, draw=secondary, line width=1.8pt, fill=secondary!16,
               text=secondary!80!black, font=\small\bfseries, align=center,
               minimum width=3.0cm, minimum height=1.0cm},
  arr/.style={-{Stealth[length=2.4mm]}, secondary!60, line width=1pt},
  note/.style={font=\scriptsize\itshape, text=secondary!80!black, align=center,
               fill=white, inner sep=2pt}
]
  \node[cbox] (c0) at (0.00,0) {Ch.\ 5\\Recurrences I};
  \node[hbox] (c1) at (3.80,0) {Ch.\ 6\\Master Theorem};
  \node[cbox] (c2) at (7.60,0) {Ch.\ 7\\Brute Force};
  \node[cbox] (c3) at (11.40,0) {Ch.\ 8\\Merge Sort};
  \node[cbox] (c4) at (15.20,0) {Ch.\ 9\\Quicksort};
  \draw[arr] (c0) -- (c1);
  \draw[arr] (c1) -- (c2);
  \draw[arr] (c2) -- (c3);
  \draw[arr] (c3) -- (c4);
  \node[note] at (1.90,-1.0) {the shortcut};
  \node[note] at (5.70,-1.0) {the baseline};
  \node[note] at (9.50,-1.0) {the pattern};
  \node[note] at (13.30,-1.0) {and its rival};
  \node[font=\scriptsize\itshape, text=accent, align=center] at (3.80,1.15)
       {solves most recurrences at sight};
\end{tikzpicture}}%
{Reading path through Part~II. Chapter 6 is the tool you will reach for most; Chapter 7 is the baseline that makes the rest worth having.}{part2map}

\newpage
